
\begin{frame}[fragile]{SQL, première partie}

Cette session présente les requêtes SQL \alert{conjonctives}, celles
qui s'expriment  sans négation ni disjonction.

\vfill

Dans cette session:

\vfill

\begin{redbullet}
  \item Forme d'un  requête SQL: variables-nuplet, conditions, construction du nuplet-résultat
  \item Requête mono-variables
  \item Requêtes multi-variables
\end{redbullet}

\vfill
\begin{blocgauche}
\alert{Ces diapositives correspondent au support en ligne disponible sur
le site \url{http://sql.bdpedia.fr/}}
\end{blocgauche}

\end{frame}


%%%%%%ù

\begin{frame}[fragile]{Variables nuplet}

SQL manipule des nuplets libres de la forme $t(a_1, a_2, \cdots, a_n)$
Nous les appellerons
des \alert{variables nuplet}.

\vfill

La quantification porte sur la variable nuplet $t$: $\exists t$ et $\forall t$.

\vfill

On désigne les attributs en les rattachant à $t$: $t.a_1$, $t.a_2$, etc.

\vfill

On  peut exprimer des comparaisons: $t.a_i = 'a'$ ou $t.a_i = t.a_j$

\end{frame}

%%%%%%%
\begin{frame}[fragile]{Requête mono-variable}

Les requêtes les plus simples utilisent une seule variable nuplet. 
Leur forme logique est:


 $$\{ t.a_1, t.a_2, \cdots, t.a_n | T(t) \land F_{cond}(t) \}$$
 
\vfill
La forme SQL:

\begin{minted}[baselinestretch=.9,
               fontsize=\small,
               framesep=2mm]{sql}
     select [distinct] t.a1, t.a2, ..., t.an
     from T as t
     where <condition>
\end{minted}

\vfill

\begin{blocgauche}
C'est un ``bloc'' avec trois clauses:
\begin{redbullet}
   \item le \texttt{from} définit la variable libre et sa \alert{portée}
   \item  le \texttt{where} définit la condition sur la variable libre
   \item le \texttt{select} (avec \texttt{distinct} optionnel) construit le nuplet-résultat
\end{redbullet}

\end{blocgauche}
\end{frame}


%%%%%%%%%%%%%%%%%%%%
\begin{frame}[fragile]{Parlons du \texttt{distinct}}

Une relation n'a pas de  doublons. Or certaines requêtes peuvent en produire:

\vfill

\begin{minted}[baselinestretch=.9,
               fontsize=\small,
               framesep=2mm]{sql}
select l.type from Logement as l
\end{minted}

\vfill

\begin{tabular}{c}
\textbf{type}\\ \hline 
Auberge\\ \hline 
Hôtel\\ \hline 
Gîte\\ \hline 
Hôtel\\ \hline 
\end{tabular}

\vfill
Le \texttt{distinct} garantit que les doublons sont éliminés.
\vfill
\begin{minted}[baselinestretch=.9,
               fontsize=\small,
               framesep=2mm]{sql}
select distinct  l.type from Logement as l
\end{minted}

\vfill

\begin{blocgauche}
Certaines requêtes ne peuvent pas produire de doublon! À approfondir.
\end{blocgauche}
\end{frame}



%%%%%%%%
\begin{frame}[fragile]{Premier exemple}

Code, nom et type des logements en Corse.

\vfill

\begin{minted}[baselinestretch=.9,
               fontsize=\small,
               framesep=2mm]{sql}
      select t.code, t.nom, t.type
      from Logement as t
      where t.lieu = 'Corse'
\end{minted}

\vfill

Correspond à la formule
$$\{ t.code, t.nom, t.type | Logement(t) \land t.lieu=\text{'Corse'} \}$$

\vfill

Forme simplifiée. 

\begin{minted}[baselinestretch=.9,
               fontsize=\small,
               framesep=2mm]{sql}
      select code, nom, type
      from Logement
      where lieu = 'Corse'
\end{minted}

\end{frame}


%%%%%%%%
\begin{frame}[fragile]{Interprétation}

La variable peut être affectée à tous les nuplets  de la table définie
par la portée.

\vfill

\alert{On garde toutes les affectations qui satisfont la condition $F_{cond}$}.

\vfill

La seule affectation correcte est surlignée ci-dessous.


\begin{center}
\begin{tabular}{c|c|c|c|c}
\textbf{code} & \textbf{nom} & \textbf{capacité} & \textbf{type} & \textbf{lieu}\\ \hline 
\rowcolor{asparagus}
pi & U Pinzutu & 10 & Gîte & Corse\\ \hline 
ta & Tabriz & 34 & Hôtel & Bretagne\\ \hline 
ca & Causses & 45 & Auberge & Cévennes\\ \hline 
ge & Génépi & 134 & Hôtel & Alpes\\ \hline 
\end{tabular}
\end{center}


\begin{blocgauche}
\alert{Trivial?} Oui, et tant mieux, car cette interprétation
fonctionne pour \alert{toutes} les requêtes.
\end{blocgauche}
\end{frame}


%%%%%%%
\begin{frame}[fragile]{Requête multi-variables}

Regardons pour deux variables: la généralisation est facile.

\vfill

Forme de la requête:

\begin{minted}[baselinestretch=.9,
               fontsize=\small,
               framesep=2mm]{sql}
select [distinct] t1.a1, ..., t1.an, t2.a1, ..., t2.am
from T1 as t1, T2 as t2
where <condition>
\end{minted}

\vfill
\begin{blocgauche}
\begin{defblock}{Interprétation}
Parmi toutes les affectations possibles des variables, 
on ne conserve que celles qui satisfont la condition exprimée par  $F_{cond}$.
\end{defblock}
\end{blocgauche}

\end{frame}


%%%%%%%
\begin{frame}[fragile]{Un exemple détaillé: logements où on peut pratiquer le ski}


Nous avons besoin de deux variables:
\begin{redbullet}
   \item la première s’affecte aux nuplets de \texttt{Activité}; 
   \item la seconde s’affecte aux nuplets de \texttt{Logement}
   \item l'attribut \texttt{codeActivité} de la première est \texttt{Ski}.
   \item \alert{les deux variables partagent le même code logement}
\end{redbullet}

\vfill

\begin{minted}[baselinestretch=.9,
               fontsize=\small,
               framesep=2mm]{sql}
select l.code, l.nom
from Logement as l, Activité as a
where l.code = a.codeLogement
and   a.codeActivité = 'Ski'
\end{minted}

\end{frame}

%%%%%%%%
\begin{frame}[fragile]{Interprétation: affectation des deux variables}

\begin{tabular}{c|c|c|c|c}
\multicolumn{5}{c}{\textbf{Logement (variable $l$)}}\\ \hline
\textbf{code} & \textbf{nom} & \textbf{capacité} & \textbf{type} & \textbf{lieu}\\ \hline 
pi & U Pinzutu & 10 & Gîte & Corse\\ \hline 
ta & Tabriz & 34 & Hôtel & Bretagne\\ \hline 
ca & Causses & 45 & Auberge & Cévennes\\ \hline 
\rowcolor{asparagus}
ge & Génépi & 134 & Hôtel & Alpes\\ \hline 
\end{tabular}

\vspace*{0.3cm}



\begin{tabular}{c|c}
\multicolumn{2}{c}{\textbf{Activité  (variable $a$)}}\\ \hline
\textbf{codeLogement} & \textbf{codeActivité}\\ \hline 
ca & Randonnée\\ \hline 
ge & Piscine\\ \hline 
\rowcolor{bananamania}
ge & Ski\\ \hline 
pi & Plongée\\ \hline 
pi & Voile\\ \hline 
\end{tabular}

\end{frame}


%%%%%%%
\begin{frame}[fragile]{Deuxième exemple: les paires de logements qui sont du même type}


Nous avons besoin de deux variables, 
\begin{redbullet}
   \item \alert{chacune ayant pour portée la table \texttt{Logement}}
   \item \alert{les deux variables partagent le même attribut \texttt{type}}
\end{redbullet}

\vfill

\begin{minted}[baselinestretch=.9,
               fontsize=\small,
               framesep=2mm]{sql}
select distinct l1.nom as nom1, l2.nom as nom2
from Logement as l1, Logement as l2
where l1.type = l2.type
\end{minted}

\vfill
\begin{blocgauche}
Soit la formule

$$\{ l_1.nom, l_2.nom | Logement(l_1) \land Logement(l_2) \land l_1.type = l_2.type \}$$
\end{blocgauche}
\end{frame}


%%%%%%%%
\begin{frame}[fragile]{Interprétation: affectation des deux variables}

\begin{tabular}{c|c|c|c|c}
\multicolumn{5}{c}{\textbf{Logement (variable $l_1$)}}\\ \hline
\textbf{code} & \textbf{nom} & \textbf{capacité} & \textbf{type} & \textbf{lieu}\\ \hline 
pi & U Pinzutu & 10 & Gîte & Corse\\ \hline 
\rowcolor{asparagus}
ta & Tabriz & 34 & Hôtel & Bretagne\\ \hline 
ca & Causses & 45 & Auberge & Cévennes\\ \hline 
ge & Génépi & 134 & Hôtel & Alpes\\ \hline 
\end{tabular}

\vfill

\begin{tabular}{c|c|c|c|c}
\multicolumn{5}{c}{\textbf{Logement (variable $l_2$)}}\\ \hline
\textbf{code} & \textbf{nom} & \textbf{capacité} & \textbf{type} & \textbf{lieu}\\ \hline 
pi & U Pinzutu & 10 & Gîte & Corse\\ \hline 
ta & Tabriz & 34 & Hôtel & Bretagne\\ \hline 
ca & Causses & 45 & Auberge & Cévennes\\ \hline 
\rowcolor{bananamania}
ge & Génépi & 134 & Hôtel & Alpes\\ \hline 
\end{tabular}

\end{frame}

%%%%%%%%
\begin{frame}[fragile]{Autre affectation possible}

\begin{tabular}{c|c|c|c|c}
\multicolumn{5}{c}{\textbf{Logement (variable $l_1$)}}\\ \hline
\textbf{code} & \textbf{nom} & \textbf{capacité} & \textbf{type} & \textbf{lieu}\\ \hline 
pi & U Pinzutu & 10 & Gîte & Corse\\ \hline 
ta & Tabriz & 34 & Hôtel & Bretagne\\ \hline 
ca & Causses & 45 & Auberge & Cévennes\\ \hline 
\rowcolor{asparagus}
ge & Génépi & 134 & Hôtel & Alpes\\ \hline 
\end{tabular}

\vfill

\begin{tabular}{c|c|c|c|c}
\multicolumn{5}{c}{\textbf{Logement (variable $l_2$)}}\\ \hline
\textbf{code} & \textbf{nom} & \textbf{capacité} & \textbf{type} & \textbf{lieu}\\ \hline 
pi & U Pinzutu & 10 & Gîte & Corse\\ \hline 
\rowcolor{bananamania}
ta & Tabriz & 34 & Hôtel & Bretagne\\ \hline 
ca & Causses & 45 & Auberge & Cévennes\\ \hline 
ge & Génépi & 134 & Hôtel & Alpes\\ \hline 
\end{tabular}

\end{frame}

%%%%%%%%
\begin{frame}[fragile]{Encore une autre (et trois autres encore possibles)}

\begin{tabular}{c|c|c|c|c}
\multicolumn{5}{c}{\textbf{Logement (variable $l_1$)}}\\ \hline
\textbf{code} & \textbf{nom} & \textbf{capacité} & \textbf{type} & \textbf{lieu}\\ \hline 
pi & U Pinzutu & 10 & Gîte & Corse\\ \hline 
ta & Tabriz & 34 & Hôtel & Bretagne\\ \hline 
\rowcolor{asparagus}
ca & Causses & 45 & Auberge & Cévennes\\ \hline 
ge & Génépi & 134 & Hôtel & Alpes\\ \hline 
\end{tabular}

\vfill

\begin{tabular}{c|c|c|c|c}
\multicolumn{5}{c}{\textbf{Logement (variable $l_2$)}}\\ \hline
\textbf{code} & \textbf{nom} & \textbf{capacité} & \textbf{type} & \textbf{lieu}\\ \hline 
pi & U Pinzutu & 10 & Gîte & Corse\\ \hline 
ta & Tabriz & 34 & Hôtel & Bretagne\\ \hline 
\rowcolor{bananamania}
ca & Causses & 45 & Auberge & Cévennes\\ \hline 
ge & Génépi & 134 & Hôtel & Alpes\\ \hline 
\end{tabular}

\end{frame}
%%%%%%%%%
\begin{frame}{À retenir}

Quelle que soit sa complexité, l'interprétation d'une requête SQL peut \alert{toujours}
se faire de la manière suivante.

\vfill

\begin{redbullet}
  \item  Chaque variable du \texttt{from} peut être affectée à tous les nuplets
       de sa portée.
  \item Le \texttt{where} définit une condition sur ces variables: seules
     les affectations satisfaisant cette condition sont conservées
    \item Le nuplet résultat est construit à partir de ces affectations
\end{redbullet}

\end{frame}

